\section{从零构建 Coding Agent}

\subsection{什么是 Coding Agent}

Coding Agent 是一种能够自主执行编程任务的 AI 系统。其核心组件包括：

\begin{itemize}
    \item \textbf{System Prompt}：定义 Agent 的行为和指令
    \item \textbf{User Prompt}：用户的自定义请求
    \item \textbf{Assistant Prompt}：LLM 的响应
    \item \textbf{Tool Set}：Agent 可调用的工具集合
\end{itemize}

\begin{importantbox}{Coding Agent 的核心循环}
一个 Coding Agent 的基本工作流程极其简单：
\begin{enumerate}
    \item 从终端读取用户输入，持续追加到对话历史
    \item 告知 LLM 可用的工具（Read\_file、List\_dir、Edit\_file 等）
    \item LLM 在适当时机请求调用工具
    \item 离线执行工具调用并返回结果
    \item 重复上述过程直到任务完成
\end{enumerate}
\end{importantbox}

\subsection{``秘密武器''：Claude Code 的设计哲学}

深入分析 Claude Code 的底层设计，可以发现几个关键的工程决策：

\begin{enumerate}
    \item \textbf{前置 context 注入}：使用小而精准的 prompt 进行 context 前置加载
    \item \textbf{System reminders 无处不在}：在 system prompt、user prompt、tool calls、tool results 中都插入 \texttt{<system-reminder>} 标签，防止模型在长对话中偏离目标
    \item \textbf{Command prefix extraction}：从用户输入中提取命令前缀，实现快捷操作
    \item \textbf{Subagent 派生}：生成子 Agent 处理特定子任务，避免单个 context 过载
\end{enumerate}

\begin{warningbox}{Context Drift 问题}
在长对话中，LLM 容易``忘记''早期的指令和约束——这被称为 context drift。Claude Code 通过在对话的各个环节反复插入 system reminders 来对抗这一问题。这是一个经过工程验证的关键技巧。
\end{warningbox}

\subsection{本章小结}

构建一个基础 Coding Agent 的技术门槛其实很低——核心就是 LLM + tool calling 的循环。但要构建一个\textbf{可靠}的 Coding Agent，需要大量的工程细节：context 管理、drift 防护、子任务分解等。

